% >>> PROPAGATION (Fable 5, Sep 1 2026): the NO-CUT reframe carried here from the long paper and the Brief --- the GRW
% kinship a shared POSITION not a shared parameter; no cut drawn; the dependency a claim about records; ceiling / ladder /
% threshold / rung retired. See the ledger, R123.
% >>> PROPAGATION (Fable 5, Aug 29 2026): the filtration clarification added at \S transit --- permanence does not separate the
% pictures by itself; it makes the placement of delta a question with an answer (the corpus-wide R6/T3 regrade). See R109.
% Determinacy as a Two-Place Relation --- the realist no-go extraction.
% Drafted August 13, 2026 (Fable 5) to the referee's drafting brief; author-verified
% per the corpus provenance convention. Companion citation policy: AethusBrief primary,
% Aethus once. Hard constraints honored: no Powerset theorem, no Nexus, no two-tag
% interferometry, no visibility formulas, no naming apparatus; "information state"
% throughout; no footnote over four sentences; Spekkens quoted at few-words grade only.
% ----------------------------------------------------------------------------
% GRADING PASS (Fable 5, August 29, 2026), propagated from the Brief's second review
% and its Opus referee response. Fenced "% >>> EDIT Sn" ... "% <<< END EDIT Sn".
%   (S1) Abstract: the no-go's summary graded --- trace-free deflation requires
%        overwrite; denied overwrite, the reparse is met horn by horn (inside a
%        persisting state: contradiction with its content; outside every state:
%        single-sample-space model). The body already said this; the abstract
%        did not.
%   (S2) \S nogo subsection title: "Deflation requires overwrite" narrowed to
%        "Trace-free deflation requires overwrite", matching the lemma it heads.
%   (S3) \S hatch, option 3: the "two bookkeepings" parsimony count marked as
%        the framework's side of a dispute, with the decoherence realist's
%        counter-count on the page.
%   (S4) \S grid, Frauchiger--Renner: a factual correction. The protocol never
%        opens the laboratories or pools the friends' reports; the superobservers
%        measure the laboratories in a basis that coherently ERASES the friends'
%        records, and it is that erasure the trigger denies for observer-grade
%        systems. The "later exit is a retrieval" clause survives only as a
%        remark about a different protocol.
%   (S5) \S limits, irretrievability: the ceiling graded as a commitment to a
%        prediction (becomes one when the construction names a rung), with the
%        two disciplines the Brief now carries (echo regime; ladder from below
%        cannot confirm).
% ============================================================================
\begin{filecontents}{realist_refs.bib}
@misc{AethusSemantics, author={Benander, Ajax}, title={Counterfactual Evaluation by Reopening Information States}, year={2026}, howpublished={Preprint}, url={https://philpapers.org/rec/BENCEB-2}}
% Entries verified August 2026 (author-side search + referee knowledge-check); the
% standing convention's final eyeball at submission still applies.
@misc{AethusQuant, author={Benander, Ajax}, title={Quantum Mechanics from Records: Measurement Structure on a Finite Path-Kernel Domain}, year={2026}, howpublished={Preprint}, url={https://philpapers.org/rec/BENTQB}}
@misc{AethusBrief, author={Benander, Ajax}, title={Aethic Reasoning in Brief: From a Counterfactual Semantics to the Quantum Measurement Problem}, year={2026}, howpublished={Preprint}, url={https://philpapers.org/rec/BENARI-3}, eprint={https://philsci-archive.pitt.edu/id/eprint/30439}, eprinttype={Alternative URL}}
@misc{Aethus, howpublished={Manuscript}, author={Benander, Ajax}, title={Aethic Reasoning: A Comprehensive Solution to the Quantum Measurement Problem}, year={2024}, month = {11}, url={https://philpapers.org/rec/BENARA-5}, eprint={https://philsci-archive.pitt.edu/id/eprint/25713}, eprinttype={Alternative URL}, note={Manuscript; the complete development of the framework, successive deposits November 2024--2026}}
@misc{AethusUnion, author={Benander, Ajax}, title={Exclusivity as Achievement: Relocating Mutual Exclusion to the Interior, and the Composition Regime That Falls Out}, year={2026}, howpublished={Preprint}, url={https://philpapers.org/rec/BENEAA-9}}
% Title read from the extraction's source (Aug 2026); closed.

@article{bell1964einstein,
  author = {Bell, John S.},
  title = {On the {Einstein}--{Podolsky}--{Rosen} paradox},
  journal = {Physics Physique Fizika},
  volume = {1}, number = {3}, pages = {195--200}, year = {1964}
}
@article{fine1982hidden,
  author = {Fine, Arthur},
  title = {Hidden variables, joint probability, and the {Bell} inequalities},
  journal = {Physical Review Letters},
  volume = {48}, number = {5}, pages = {291--295}, year = {1982}
}
@article{pitowsky1994george,
  author = {Pitowsky, Itamar},
  title = {George {Boole}'s `conditions of possible experience' and the quantum puzzle},
  journal = {British Journal for the Philosophy of Science},
  volume = {45}, number = {1}, pages = {95--125}, year = {1994}
}
@article{kochen1967problem,
  author = {Kochen, Simon and Specker, Ernst P.},
  title = {The problem of hidden variables in quantum mechanics},
  journal = {Journal of Mathematics and Mechanics},
  volume = {17}, number = {1}, pages = {59--87}, year = {1967}
}
@article{spekkens2007evidence,
  author = {Spekkens, Robert W.},
  title = {Evidence for the epistemic view of quantum states: A toy theory},
  journal = {Physical Review A},
  volume = {75}, pages = {032110}, year = {2007}
}
@misc{spekkens2018quantum,
    author = "Spekkens, Robert",
    title = "{Robert Spekkens Public Lecture: The Riddle of the Quantum Sphinx}",
    howpublished = "YouTube",
    year = "2018",
    month = "February",
    day = "8",
    note = "Perimeter Institute Public Lecture Series",
    url = "https://www.youtube.com/watch?v=jvZ5m3UMer4"
}
% Second source in the transcript on file --- uncited in this paper; canonical entry kept
% here for the record.
@misc{spekkens2024escapedsapiens,
  author       = {Spekkens, Robert W.},
  title        = {What is Reality Made Of? {Unscrambling} the Quantum Omelet},
  howpublished = {Interview by Shane Farnsworth, The Escaped Sapiens Podcast, Episode 74},
  year         = {2024},
  month        = nov,
  note         = {Video available at \url{https://youtu.be/J2ZIRkfrFlI}},
  url          = {https://youtu.be/J2ZIRkfrFlI}
}
@article{rovelli1996relational,
  author = {Rovelli, Carlo},
  title = {Relational quantum mechanics},
  journal = {International Journal of Theoretical Physics},
  volume = {35}, number = {8}, pages = {1637--1678}, year = {1996}
}
@article{adlamrovelli2022,
  author = {Adlam, Emily and Rovelli, Carlo},
  title = {Information is physical: cross-perspective links in relational quantum mechanics},
  journal = {Philosophy of Physics},
  volume = {1}, number = {1}, pages = {4}, year = {2023},
  note = {arXiv:2203.13342 (2022); AUTHOR TO VERIFY article number at submission}
}
@article{fuchs2014introduction,
  author = {Fuchs, Christopher A. and Mermin, N. David and Schack, R{\"u}diger},
  title = {An introduction to {QBism} with an application to the locality of quantum mechanics},
  journal = {American Journal of Physics},
  volume = {82}, number = {8}, pages = {749--754}, year = {2014}
}
@incollection{fine2005tense,
  author = {Fine, Kit},
  title = {Tense and reality},
  booktitle = {Modality and Tense: Philosophical Papers},
  publisher = {Oxford University Press}, address = {Oxford},
  pages = {261--320}, year = {2005}
}
@article{wood_spekkens_2015,
  author = {Wood, Christopher J. and Spekkens, Robert W.},
  title = {The lesson of causal discovery algorithms for quantum correlations: causal explanations of {Bell}-inequality violations require fine-tuning},
  journal = {New Journal of Physics},
  volume = {17}, pages = {033002}, year = {2015}
}
@article{peres1978unperformed,
  author = {Peres, Asher},
  title = {Unperformed experiments have no results},
  journal = {American Journal of Physics},
  volume = {46}, number = {7}, pages = {745--747}, year = {1978}
}
% Confirmed published (Aug 2026); arXiv kept as eprint.
@article{zwirn2023convivial,
  author = {Zwirn, Herv{\'e}},
  title = {Everett's interpretation and convivial solipsism},
  journal = {Quantum Reports},
  volume = {5}, number = {1}, pages = {267--281}, year = {2023},
  doi = {10.3390/quantum5010018},
  eprint = {2301.07532}
}
% Confirmed by author-side search (Aug 2026).
@misc{ics2026dissolution,
  author = {Sendall, Jonathon},
  title = {Interaction-conditional semantics and the dissolution of quantum paradoxes},
  howpublished = {arXiv:2601.18810 [quant-ph]},
  year = {2026}
}
@article{auffeves2016contexts,
  author = {Auff{\`e}ves, Alexia and Grangier, Philippe},
  title = {Contexts, systems and modalities: A new ontology for quantum mechanics},
  journal = {Foundations of Physics},
  volume = {46}, number = {2}, pages = {121--137}, year = {2016}
}
@article{healey2012quantum,
  author = {Healey, Richard},
  title = {Quantum theory: A pragmatist view},
  journal = {British Journal for the Philosophy of Science},
  volume = {63}, number = {4}, pages = {729--771}, year = {2012}
}
@article{bohm1952suggested,
  author = {Bohm, David},
  title = {A suggested interpretation of the quantum theory in terms of ``hidden'' variables. {I}},
  journal = {Physical Review},
  volume = {85}, number = {2}, pages = {166--179}, year = {1952}
}
@article{everett1957relative,
  author = {Everett, Hugh},
  title = {``Relative state'' formulation of quantum mechanics},
  journal = {Reviews of Modern Physics},
  volume = {29}, number = {3}, pages = {454--462}, year = {1957}
}
@article{ghirardi1986unified,
  author = {Ghirardi, GianCarlo and Rimini, Alberto and Weber, Tullio},
  title = {Unified dynamics for microscopic and macroscopic systems},
  journal = {Physical Review D},
  volume = {34}, number = {2}, pages = {470--491}, year = {1986}
}
@article{frauchiger2018quantum,
  author = {Frauchiger, Daniela and Renner, Renato},
  title = {Quantum theory cannot consistently describe the use of itself},
  journal = {Nature Communications},
  volume = {9}, pages = {3711}, year = {2018}
}
@article{pusey2012reality,
  author = {Pusey, Matthew F. and Barrett, Jonathan and Rudolph, Terry},
  title = {On the reality of the quantum state},
  journal = {Nature Physics},
  volume = {8}, number = {6}, pages = {475--478}, year = {2012}
}
@article{brukner2018nogo,
  author = {Brukner, {\v{C}}aslav},
  title = {A no-go theorem for observer-independent facts},
  journal = {Entropy},
  volume = {20}, number = {5}, pages = {350}, year = {2018}
}
@article{bong2020strong,
  author = {Bong, Kok-Wei and Utreras-Alarc{\'o}n, An{\'i}bal and Ghafari, Farzad and Liang, Yeong-Cherng and Tischler, Nora and Cavalcanti, Eric G. and Pryde, Geoff J. and Wiseman, Howard M.},
  title = {A strong no-go theorem on the {Wigner}'s friend paradox},
  journal = {Nature Physics},
  volume = {16}, number = {12}, pages = {1199--1204}, year = {2020}
}
@book{thooft2016cellular,
  author = {'t Hooft, Gerard},
  title = {The Cellular Automaton Interpretation of Quantum Mechanics},
  series = {Fundamental Theories of Physics}, volume = {185},
  publisher = {Springer}, address = {Cham}, year = {2016}
}
@article{hossenfelder2020rethinking,
  author = {Hossenfelder, Sabine and Palmer, Tim},
  title = {Rethinking superdeterminism},
  journal = {Frontiers in Physics},
  volume = {8}, pages = {139}, year = {2020}
}
@book{griffiths2002consistent,
  author = {Griffiths, Robert B.},
  title = {Consistent Quantum Theory},
  publisher = {Cambridge University Press}, address = {Cambridge}, year = {2002}
}
@article{leifer2014psi,
  author = {Leifer, Matthew S.},
  title = {Is the quantum state real? {An} extended review of $\psi$-ontology theorems},
  journal = {Quanta},
  volume = {3}, number = {1}, pages = {67--155}, year = {2014}
}
% Archive-edition route (frozen), per recommendation; first published 2002, rev. 2017.
@misc{lombardi2021modal,
  author = {Lombardi, Olimpia and Dieks, Dennis},
  title = {Modal interpretations of quantum mechanics},
  howpublished = {The Stanford Encyclopedia of Philosophy (Winter 2021 Edition), E. N. Zalta (ed.)},
  url = {https://plato.stanford.edu/archives/win2021/entries/qm-modal/},
  year = {2021}
}
% Confirmed: the 1976 second edition is the standard scholarly citation here.
@book{despagnat1976conceptual,
  author = {d'Espagnat, Bernard},
  title = {Conceptual Foundations of Quantum Mechanics},
  edition = {2}, publisher = {W. A. Benjamin}, address = {Reading, MA}, year = {1976}
}

@misc{GoldenScience2022,
  author       = {Benander, Ajax},
  title        = {Golden Science},
  year         = {2022},
  month        = {November},
  howpublished = {Self-hosted manuscript},
  url          = {https://aethus.org/aethic-code/files/GoldenScience_Nov2022.pdf},
  note         = {Title page August 2022, the date of the first version of the source; this version created November 2022. Not deposited to a
                  preprint server; cited as a dated developmental record rather than
                  as a source of results.}
}

@article{elitzur1993quantum,
  author  = {Elitzur, Avshalom C. and Vaidman, Lev},
  title   = {Quantum Mechanical Interaction-Free Measurements},
  journal = {Foundations of Physics},
  volume  = {23},
  number  = {7},
  pages   = {987--997},
  year    = {1993},
  doi     = {10.1007/BF00736012}
}
\end{filecontents}

\documentclass[11pt]{article}
\usepackage[margin=1in]{geometry}
\usepackage{amsmath, amssymb, amsthm}
\usepackage[hidelinks]{hyperref}
\usepackage{microtype}
\usepackage{booktabs}
\usepackage[style=numeric, sorting=none]{biblatex}
\addbibresource{refs.bib}
\addbibresource{realist_refs.bib}

\newtheorem{theorem}{Theorem}
\newtheorem{lemma}{Lemma}
\newtheorem{corollary}{Corollary}
\newtheorem{definition}{Definition}
\newtheorem{proposition}{Proposition}

\title{Determinacy as a Two-Place Relation:\\ The Category Mistake in Quantum Interpretation, and the Paradoxes Re-indexed}
\author{Ajax Benander\thanks{This paper extracts the realist position from the framework of Aethic reasoning; the compressed overview \cite{AethusBrief} is the primary companion throughout, with the full development in \cite{Aethus}. The exposition was drafted in collaboration with an AI system (Claude, Anthropic) under the corpus's standing provenance conventions, and every argument has been verified by, and is owned by, the author. Correspondence: akb9408@rit.edu; \url{https://aethus.org}.}}
\date{First drafted August 2026\\ {\small This version: October 2026 --- draft}}

\begin{document}
\maketitle

\begin{abstract}
\noindent We propose that determinacy is a two-place relation between an attribute and an information state, rather than a one-place property that facts have or lack, and examine what a realist about quantum mechanics gains and loses by adopting it. Two results of limited scope carry the formal part. An internal preservation lemma: a completion that retains the framework's constructive openness of an entry cannot also assign that entry a value at the same index. A common-space corollary: the class of noncontextual, measurement-independent completions, whose image is a single sample space, is empty wherever the observed statistics violate the Bell inequalities, by Fine's theorem. Neither result selects information-state indexing from the alternatives that also reject that completion class; the case for the indexing is comparative. One criterion then carries the applications, for a detector and a laboratory alike: a candidate coherence at a stated boundary is excluded exactly when every admissible refinement of it retains a descendant stating a record, accessible outside that boundary, that distinguishes the alternatives held coherent. For the observer case the criterion's premise, that an observer-grade laboratory always retains such a descendant, is a physical commitment on which the framework departs from universal unitarity; it is stated as a premise, and what it excludes is a coherent state assignment, not an apparatus operation.
\end{abstract}

\section{Introduction: the realist's dilemma, and a shared diagnosis}\label{sec:intro}

The realist's dilemma in quantum mechanics can be stated in a sentence. If unperformed measurements have determinate results, and those results are the values of variables on one common sample space, the same whichever compatible measurements accompany them and distributed independently of the settings chosen, then the joint statistics they would have to exhibit are constrained by the Bell inequalities, and the world violates them \cite{bell1964einstein}; if unperformed measurements have no results at all \cite{peres1978unperformed}, then claims about what an experiment \emph{would have shown} need some other truth-maker than a determinate outcome, or they are not truth-apt and the realism has quietly become an instrumentalism about everything not actually done. That is a demand for a realist semantics of the unperformed experiment, and not a dilemma every interpretation must fail: probabilistic and contextual counterfactuals meet it in their own ways, the scope discussion below admits both, and the paper's own answer, a modal structure without a determinate outcome, is one more way of meeting it. Each interpretation on the standard menu takes one horn or the other, at a cost its critics can name.

\textit{The claim.} This paper claims both halves: enough counterfactual structure to say what unperformed experiments would have shown, and no hidden valuation for the Bell-type theorems to kill. The paper's burden is to show what four commitments stated in advance entail, what an established obstruction adds to them, and what the resulting position explains. The commitments are not themselves derived here; they are laid out, their costs named, and the arguments run on them.

\noindent One calibration belongs here rather than at the end, because the title states a thesis and the paper's results are narrower. Three claims are made, of different status, and the paper keeps them apart throughout.
\begin{center}\small
\begin{tabular}{p{0.60\textwidth} p{0.32\textwidth}}
\toprule
\textbf{Claim} & \textbf{Status} \\
\midrule
(i) Constructive openness cannot be preserved unchanged while a selecting value is added at the same index. & Internal preservation result (\S\ref{sec:transit}). \\
(ii) The specified common-space completion, noncontextual and measurement-independent, cannot reproduce Bell-violating statistics. & Application of an established external obstruction, Fine's (\S\ref{sec:samplespace}). \\
(iii) Indexed constructive openness explains the relevant cases more coherently, or more economically, than the alternatives. & Philosophical argument requiring comparison (\S\S\ref{sec:grid}, \ref{sec:placement}). \\
\bottomrule
\end{tabular}
\end{center}
\noindent Claim (i) is a result about the framework's own objects, and a critic who disputes the constructive reading of openness is not refuted by it. Claim (ii) excludes a completion class; it does not establish that determinacy must take an information state as its second argument. Claim (iii) is where the paper's distinctive contribution lies, and it inherits no strength from the first two. The category-mistake diagnosis---that the one-place reading of determinacy misapplies a predicate that takes two places---is the thesis; (i) and (ii) are the arguments for it that can be made rigorous, and the applications of \S\ref{sec:grid} are graded as consequences of (i) and (ii) under named further premises, not as independent proofs of the thesis. One dependency is recorded now because later sections rely on it: (ii) does not depend on the permanence of predecessor states, and (i) arises from the constructive-openness commitment alone, at a single index, without any historical update. The permanence commitment, L4 below, serves the ontology and the historical evaluation of counterfactuals; it is not indispensable to either result.



\textit{The shared diagnosis.} The argument is organized around a diagnosis this paper shares with Spekkens: that the interpretation debate is stalled on a category mistake, an error not about some entity but about what kind of entity it is \cite{spekkens2018quantum}. His framing has four parts, used here as motivation and not as evidence. A completed theory is recognizable by the mutual support of its parts, as a solved puzzle is recognized by checking that all of its constraints are satisfied at once. A century without consensus on the interpretation of quantum theory motivates re-examination of early assumptions; it demonstrates nothing by itself, and in particular it does not show that any existing interpretation fails. The incumbents can then be read as internally coherent work built on an early entry that was never argued for. And a historical calibration, the ideographic reading of Egyptian hieroglyphs, shows what the correction of such an entry looks like: not an answer to the old questions but their replacement. The diagnosis is made here within the proposed conception of determinacy; the internal argument cannot by itself move a reader who rejects that conception, so the case made to such a reader is comparative, claim (iii) above (\S\ref{sec:placement}, \S\ref{sec:limits}).

The diagnosis, taken as correct, poses a single question: \emph{which} category was mistaken? This paper's answer is that determinacy has been treated as a \emph{one-place property} --- something spacetime-located facts simply have or lack --- when it is a \emph{two-place relation} between an attribute and an information state. The information states in question are not anyone's credences: they are logical objects, individuated by stated content, and what stands determinate relative to them is a fact about the world's structure. The position that results is a realism. \S\ref{sec:positive} states the proposal positively, with an example and the reasons for it; \S\S\ref{sec:deflation}--\ref{sec:hatch} establish claims (i) and (ii) against the strongest deflationary reading and compare the selective alternatives; \S\ref{sec:grid} runs the applications; \S\ref{sec:objections} answers the objections the proposal most invites; and \S\ref{sec:placement} places it beside its neighbors, Spekkens's own program among them, which answers the same question differently.

Two constraints govern the use of the non-consensus heuristic. It motivates looking for a wrong early assumption rather than for another downstream repair, and it says nothing about a proposal the field has not yet weighed; and the completed-theory criterion awards no verdict to this paper, since completion is recognized by a field and not claimed by an author. The paper's part is therefore to display the fit --- one correction across the paradox set, with the case-specific inputs each application consumes recorded beside it (\S\ref{sec:grid}) --- and to name what remains open (\S\ref{sec:limits}).

One scope declaration completes the introduction, and it is load-bearing. This paper argues that a semantics --- a way of assigning truth conditions to claims about the determinate, the open, and the counterfactual --- is committed to nonclassical structure of a specific kind. Whether that structure is \emph{physically instantiated} is a separate question, developed in the physics companion and surveyed in the overview \cite{AethusQuant, AethusBrief}; nothing below computes an interference pattern or classifies a laboratory boundary. The paper is philosophy of physics in the strict sense: its subject is what the realist may coherently say, its method is to prove that one specified thing she might like to say cannot be said, and its evidence is what saying the other thing explains. The one thing is narrow: the specified noncontextual, measurement-independent completion cannot preserve the required statistics, and this obstruction does not uniquely select information-state indexing from all remaining interpretations.

\section{The proposal, stated positively}\label{sec:positive}

\noindent The thesis is easiest to assess on its own, before any rival is criticized. It has three layers, kept apart below: what it asks a reader to believe, how the framework expresses it formally, and why a reader should believe it.

\textit{What it asks a reader to believe.} Whether an attribute is determinate is a fact about the attribute together with an information state, not about the attribute alone. Relative to a given state an attribute is either stated, and then determinate with its value, or open; there is no further, state-independent fact about whether it is determinate. The information states are logical objects, individuated by what they state, and what stands determinate relative to them is a fact about the world's structure, not about anyone's belief.

\textit{One example.} A coin was tossed yesterday and the result sealed in an envelope no one has opened. On the one-place view the face is determinate, it landed one way, and the envelope hides which; whoever has not read it is ignorant of a determinate value. On the two-place view, relative to the state of anyone who has not read the envelope the face is open, not determinate-but-unknown, and relative to the state that includes the reading it is determinate. Three things change. ``The face is determinate'' without an index becomes shorthand for determinacy relative to a contextually supplied state, as ``simultaneous'' is in relativity. The open entry conceals no value, there being no hidden fact at that index for it to conceal. And nothing conflicts: the reader's later state refines the earlier one by stating what it left open, while the earlier state persists as what was settled then. In the coin case the two views issue the same statistics, and the paper does not claim otherwise; the difference becomes physical where openness takes its other logical form.

\textit{What makes the blank constructive.} An open entry is not an omitted one. It is a description of a specific kind: that the attribute's alternatives stand open, with the weights the state assigns them, and, in the \emph{agreeing} form, that they stand open together, with no stated exclusivity between them. Two states that carry the same elementary value assignments but differ in their regime declarations, whether an open attribute's alternatives stand in the agreeing or the exclusive form, are different states, and their consequences can differ: an agreeing superposition of which-path alternatives carries the interference that an exclusive one, a mixture over determinate paths, does not. That is the sense in which openness is content rather than a gap: some things are true of a state because an attribute is open in a given form, and replacing the agreeing form by the exclusive form over the same values, the preparation by the corresponding mixture, does not preserve them. Whether some richer completion, contextual or nonlocal, preserves them is a separate question, and \S\ref{sec:scope} leaves it open.

\textit{How the framework expresses it.} An information state is a deductively closed partial assignment; determinacy relative to it is statedness; refinement adds stated content; and validity is governed by the third postulate, which asks of a candidate state whether some declared refinement of it is safe, all of that refinement's declared refinements being valid. The ledger of the next section fixes these as four commitments, L1 to L4.

\textit{Why accept it.} The reason is what the proposal delivers, at stated costs. It supplies indexed modal verdicts without a standing total valuation for the Bell-type theorems to exclude: in the scenario fragment, an alternative supported by some but not all visible pockets yields \emph{might} without \emph{would}, and numerical grades require a weighting with a defined normalization (Section~\ref{sec:grid}). It gives one criterion for when a record excludes a coherence, the same for a detector and a laboratory, with what each application must additionally establish stated beside it. And it treats the unread past and the unsettled future alike, openness carrying no tense. The costs are as definite: determinacy loses its one-place form, earlier information states are retained as objects, and the results of Sections~\ref{sec:nogo} and~\ref{sec:hatch} exclude only specified rivals. A reader who weighs the yield against the costs and prefers the costs is not refuted; Section~\ref{sec:objections} answers the objections such a reader is most likely to press.

\section{The ledger}\label{sec:ledger}

The argument consumes four commitments of its own, together with the standard premises of the theorems it invokes---noncontextuality where Kochen--Specker enters, measurement independence, and the quantum statistics themselves---and the four are granted here as an explicit ledger, since a conditional argument with a stated antecedent is auditable in a way an ambient worldview is not. An opponent of everything below can locate the exact entry she denies, and \S\ref{sec:nogo} will show that, short of denying one of those standard premises, locating it is the only exit. Contextual objective values, which deny noncontextuality, remain a live exit of that other kind.

\emph{L1, indexation.} Determinacy is a two-place relation: an attribute is determinate, or open, \emph{relative to an information state}---an \emph{Aethus}, plural \emph{Aethae}, in the framework's own name for the object, which is used below wherever it is the framework's specific object that matters and abbreviated to ``state'' otherwise. An Aethus, throughout, is a logical object --- a deductively closed body of stated content, individuated by what it states --- not a mental state, not a credence function, and not anyone's awareness, and equally not a property attributed to arbitrary physical systems: the working example is the standing content of a laboratory's record \emph{as it stands to those who hold it}, and no entropy or channel notion appears anywhere in the specification. Nothing below turns on where such states arise, a question this paper leaves open and its companion declines to prejudge. The one-place locution ``the attribute is determinate,'' where it survives at all, is shorthand for determinacy relative to a contextually supplied state, exactly as ``simultaneous'' survives in relativity as shorthand for simultaneity relative to a frame.

\emph{L2, tense-free openness.} An open entry --- a \emph{blank} --- is not an unknown value with a hidden tense. It is valueless relative to its state, and it resolves only in the moment of learning, resolution by deduction included; before that moment the alternatives stand genuinely open in the Markovian sense that either admissible resolution remains available while the entry is blank, however ancient the event concerned. Distance in time is not what fixes an entry; entailment from the state's content is. And the indexing holds in the present continuous, which is its sharpest form: even as one stares at a fully determinate scene, relative to one's state of thirty seconds ago the matter remains open --- determinacy gained is always indexed to the state that gained it, a permanence L4 is about to make structural. The third postulate the applications consume is the physics companion's $P_3$, stated here once, in its safety form: a valid state has a declared descendant all of whose declared descendants are valid, $\mathcal{V}(A) \Rightarrow \exists B \le A\ \forall C \le B\ \mathcal{V}(C)$, declared descent $\le$ being reflexive and transitive. The implication alone does not fix the predicate, since the universal predicate satisfies it; the validity predicate is read base-generated, and the construction is this. Let $U$ be the domain of states and $\mathrm{Base} \subseteq U$ the base contradictions, the states that jointly assert incompatible content. For $S \subseteq U$ put
\[
\Phi(S) = \{A \in U \setminus \mathrm{Base} : \exists B \le A\ \forall C \le B,\ C \in S\},
\]
which is monotone in $S$, and let $\mathcal{V} = \nu\Phi$ be its greatest fixed point, the union of all $S$ with $S \subseteq \Phi(S)$. Then $\mathcal{V}$ excludes the base, satisfies the safety clause, and contains every predicate that does both; the applications import it in exactly that form. The permanence of a predecessor as a semantic object is L4's commitment and not a consequence of this construction, and neither it nor the physical irreversibility of records may be substituted for the quantifier rule.

\begin{lemma}[Conditional record exclusion]\label{lem:record-exclusion}
Let $A_{\mathrm{coh}}$ be a candidate state asserting a specified maintained coherence, and let the relevant base contradictions be those produced when that coherence and an incompatible separating record are jointly stated. If every declared refinement of $A_{\mathrm{coh}}$ has a further declared descendant exhibiting that incompatibility,
\[
\forall B \le A_{\mathrm{coh}}\ \exists C \le B:\ C \in \mathrm{Base},
\]
then $A_{\mathrm{coh}}$ is invalid under the base-generated $P_3$ rule.
\end{lemma}
\begin{proof}
Since $\mathcal{V} = \Phi(\mathcal{V})$, validity of $A_{\mathrm{coh}}$ requires a safe refinement, a $B \le A_{\mathrm{coh}}$ all of whose declared descendants lie in $\mathcal{V} \subseteq U \setminus \mathrm{Base}$; by hypothesis every $B \le A_{\mathrm{coh}}$ has a descendant in the base, so no safe refinement exists.
\end{proof}

\noindent\emph{Converse.} For a candidate that is not itself a base contradiction the condition is also necessary. If some $B_0 \le A_{\mathrm{coh}}$ has no descendant in the base, then $S = \mathcal{V} \cup \{D : D \le B_0\}$ satisfies $S \subseteq \Phi(S)$: each $D \le B_0$ lies outside the base, and $D \le D$ by reflexivity has every descendant in $S$ by transitivity. So $S \subseteq \nu\Phi$, and $A_{\mathrm{coh}} \in \Phi(\mathcal{V}) = \mathcal{V}$ by the refinement $B_0$. A candidate outside the base is therefore excluded exactly when every declared refinement of it has a base-reaching descendant, and only the reflexivity and transitivity of declared descent are used.

\noindent The lemma is the reusable framework result; the detector and the friend are two proposed physical instantiations of its premise. ``Declared refinement'' names a \emph{proper child} of $A$, written $B \le A$: a component carried, in some admissible declared decomposition of $A$, at a weight that is nonzero and not negligible relative to the remainder's, so that a valid component carried at relatively negligible weight is no proper child and supplies no validity floor. The relation is transitive under the algebra companion's admissibility and coefficient-composition assumptions, so a proper child at any depth is a proper child, and ``two generations'' names the third postulate's two nested quantifiers, each over the full transitive relation, not a two-step horizon; the coefficient conventions are recorded in Appendix~\ref{app:properchild}. The criterion does not restrict which physical degrees of freedom may carry a separating record. Any proposed physical realization must supply its declared refinements and record incompatibilities consistently with the candidate configuration, and where those structures depend on the configuration itself, their joint consistency is an additional construction problem that the lemma does not address.

\emph{L3, the constructive reading of blanks.} A blank is not an absence of description but a description of a specific kind: a determinate \emph{entry} whose content is the openness of its alternatives, carried with whatever graded weight the state assigns them. Openness comes in two logical forms --- with a stated exclusivity clause, or without one --- and the difference, developed at full length in the dedicated extraction \cite{AethusUnion} and surveyed in the companion \cite{AethusBrief}, will matter exactly once below. Two states with the same elementary value assignments that differ in their regime declarations or other structural content are different states, describing differently --- openness is content, not a gap where content should be.

\emph{L4, transit, not overwrite.} Learning produces a successor state in which the resolved entry stands stated; it does not revise the predecessor in place. Prior states are permanent objects: the description ``open, with these alternatives at these weights'' is never destroyed, only succeeded. Nothing in what follows is owed to psychology here --- the commitment concerns the semantics' objects: it is a commitment to retaining indices so that evaluation ``from what was known then'' has a state to be evaluated at. The reopening operator itself is synchronic and may construct states never historically occupied, as the companion semantics defines it, so the operator alone does not force historical-state permanence; what forces the retention is the demand for historical evaluation.

That is the whole ledger. One further piece of machinery is imported rather than granted, because it is developed and defended elsewhere and this paper's contribution does not lie in it: the counterfactual semantics of the companion extraction \cite{AethusSemantics}, on which evaluating \emph{if $X$ had happened, $Y$ would have happened} never leaves the evaluator --- one reopens the antecedent's entry in one's own information state at a declared scope, restates it, re-prunes, and reads the verdict off what survives, with \emph{would}, \emph{might}, and \emph{the probability that} a single operator read at three grains. The import matters for one reason only: it is what gives the framework's counterfactuals literal truth conditions, so that Peres's dictum returns below with the modality distributed rather than refused --- where the alternatives are open pockets of an exclusive residue, the \emph{would} false whenever some visible pocket is a counterexample to it, the \emph{mights} true, and the grade defined where a weighting is, agreeing superpositions falling under the full framework's separate clause \cite{peres1978unperformed}.

\section{The category mistake, named}\label{sec:mistake}

The diagnostic claim of this section, made within the proposed conception of determinacy and offered to a reader outside it only through the comparison of \S\ref{sec:placement}, is that one-place determinacy occupies, in the quantum interpretation debate, exactly the structural position the ideographic assumption occupied in Egyptology: it is presupposed rather than argued; it renders a family of central questions unanswerable \emph{because malformed}; and the incumbent programs are what internally coherent work built on it looks like --- mutually incompatible with one another, and none cohering outward, which is the third movement's signature.

Presupposition first. A large and historically central family of interpretations retains or reconstructs more one-place factual structure than the present account permits, and none of that family argues \emph{for} the one-place reading; each consumes it. The claim is about that family, not about every position on the menu---the relational and contextual programs met in \S\ref{sec:placement} are exactly the ones that do not consume it whole. The hidden-variable program consumes it in its very statement, values possessed and awaiting discovery. The collapse theorist consumes it in the demand that superposition \emph{become} one-place definiteness at some point, and then owes the dynamics of the becoming \cite{ghirardi1986unified}. The Copenhagen tradition, which comes nearest to relativizing, manages the pressure with complementarity's silence rather than with a revised category: what may not be jointly said is fenced, but what saying \emph{is}, the predication of one-place facts, is never touched. The Everettian case needs a different statement, since the relative-state formulation is itself relational \cite{everett1957relative}: contemporary Everettian accounts recover definite outcomes within branches, and the disagreement with the present proposal is not about the number of argument places, ``within its world'' already being a relativization, but about what the second relatum is and what it costs. A branch is a component of the universal state individuated dynamically, by a decoherence-selected decomposition, and the recovered determinacy therefore requires the branching structure, its preferred decomposition and its reconstruction of probability, as additional ontology; an information state is individuated by stated content, carries its openness as content, and requires none of that. An assumption that a family of rivals consumes and none defends is what an early wrong entry looks like from inside the work built on it.

Malformation second. The questions that have organized a century of debate---\emph{what is the system doing between measurements}, \emph{when does the collapse occur}, \emph{along which foliation}---each consume the one-place slot, asking after a determinate fact, a doing, a moment, a slice, on the presupposition that there is one to ask after. Under the two-place correction they receive replacements rather than answers: which attributes stand open relative to which states, and at what weights; relative to which state an entry stands resolved, an indexical answer rather than a chronological one; and, for the foliation, nothing happening on a spacelike slice under re-indexing, so that no preferred slicing is owed (\S\ref{sec:grid}).

Coherent-work-on-a-wrong-entry, third. Read with the two-place candidate in hand, the standard menu resolves into internally consistent ways of protecting the one-place entry: by election of a preferred attribute family, at the cost of a nonlocal guidance dynamics \cite{bohm1952suggested}; by installment, at the cost of new dynamics with experimental tails \cite{ghirardi1986unified}; by branching, at the cost just stated; and by retreat, the quantum state reassigned to the knowledge side of a line whose worldly side remains classically conceived \cite{spekkens2007evidence}. It is not shown here that any of them fails: the non-consensus heuristic carries no verdict, and the claim against the incumbents is the comparative one, claim (iii), argued in \S\ref{sec:placement}.

One more clause completes the diagnosis, an etiology. In the everyday regime records are copied promiscuously into walls, air, instruments and memories, and every information state that has acquired them agrees with every other on the attributes they resolve; when the second relatum does not vary among the states that matter, a two-place relation is observationally indistinguishable from a one-place property, and the one-place shorthand is correct, a theorem of the record-saturated regime. The mistake was the promotion of that shorthand into the regime where states disagree, and its invisibility follows: a convention that holds as a theorem where one lives does not feel like a convention, and the first arena able to test it, the laboratory regime where record supply can be controlled, arrived three centuries after the grammar set.

The corrected entry is L1, already on the ledger. What follows is its defense: first claims (i) and (ii), against the strongest deflationary reading, then the applications.

\section{The deflationary challenge, at full strength}\label{sec:deflation}

The objection that would reduce everything above to notation is stated in its strongest form, because what \S\ref{sec:nogo} establishes is only as strong as the reading it excludes.

The deflationist speaks as follows. \emph{Your `blanks' are unknowns, and every unknown has a determinate value behind it. Your two-place `determinacy relative to a state' is one-place determinacy plus a bookkeeping of who has learned what. Apply the translation scheme $\delta$: wherever a state carries `open between $x_1, \dots, x_n$ at weights $w_1, \dots, w_n$', read `the value is some particular $x_i$, unknown to that state, with the weights expressing its ignorance'. Every inference you license survives $\delta$, every verdict your semantics issues $\delta$ reissues, and nothing observable distinguishes the images; your relativization is a coordinate transform on the classical picture. As for your permanence of prior states, the predecessor persists as a record of what was then unknown, which is what my reading predicts: memory of ignorance, not evidence of ontological openness.}

If this succeeds, the paper is a notation paper. The next section argues that it cannot succeed: the translation scheme $\delta$, read as a noncontextual, measurement-independent completion, has no action that preserves the ledger's constructive openness, and the image it would produce if it acted anywhere is excluded by the quantum statistics.

\section{The no-go}\label{sec:nogo}

Two results of different status, claims (i) and (ii) of the introduction, and one paragraph on their relation.

\subsection{The internal preservation lemma}\label{sec:transit}

\begin{lemma}[Internal preservation]\label{lem:preservation}
Grant L3. Let $I$ be a state carrying, for an attribute $q$, the constructive entry \emph{open between $x_1, \dots, x_n$ at weights $w_1, \dots, w_n$} in the agreeing form, with no exclusivity clause. A completion that retains that entry unchanged and adds, at the same index $I$, the value $q = x_k$ as holding all along produces a state asserting, of one attribute at one index, both that its alternatives stand open without exclusivity and that one of them obtains. The two entries are jointly invalid, so no such completion is a valid state.
\end{lemma}
\begin{proof}
By L3 the openness is stated content and not the absence of it: the agreeing entry states that the alternatives stand open together. A stated selecting value at the same index states that one alternative obtains and the others do not, the exclusivity the agreeing entry denies; the pair is a base contradiction by the regime rules, and a state containing it is invalid by the base-generated $\mathcal{V}$ of \S\ref{sec:ledger}.
\end{proof}

\noindent Three remarks fix the lemma's scope. It is stated at a single index and uses neither L4 nor any historical update; what L4 adds is that the historical instance, whether the morning state's blank was ignorance, remains a question about an existing object rather than one dissolved at resolution. It does not touch a deflationist who translates constructive openness into epistemic openness rather than retaining it, which is the strongest reading's own proposal; for that translation the result is the corollary below. And its force is proportional to L3: a critic who declines the constructive reading of openness is not refuted by it. One further remark concerns permanence, since the deflationist's reply above rests on it. Where an earlier description is not retained there is nothing for a reparse to act on, which is why the deflation has felt available in classical settings; a classical model that retains its earlier descriptions, a filtration keeping its past $\sigma$-algebras, is not thereby refuted, and what separates the pictures is not permanence but the regime distinction the lemma consumes: agreeing openness, whose alternatives are co-present and can interfere, against the exclusive superposition that ignorance is.

\subsection{The common-space completion corollary}\label{sec:samplespace}

\begin{definition}[Admissible deflation]\label{def:deflation}
Let $\mathcal{Q}$ be the attributes, with value spaces $V_q$, let $\mathcal{C}$ be the measurement contexts, each a set of jointly measurable attributes, with observed outcome distributions $P(\,\cdot \mid c)$ on $\prod_{q \in c} V_q$, and let $S$ be the information states, each assigning its open attributes weights. A \emph{deflation} of the framework's open entries consists of
\begin{itemize}
\item[\textbf{D1}] (\emph{Hidden sample space}) a measurable space $\Lambda$ with a probability measure $\mu$, and for every attribute $q$ a measurable value $v_q : \Lambda \to V_q$, so that every open attribute at every state receives a determinate value $\delta(s,q) = v_q(\lambda)$ at each point, with the state's weights read as a conditional measure $\mu_s$ on $\Lambda$ expressing ignorance of the point;
\item[\textbf{D2}] (\emph{Noncontextuality}) for each context $c$ and $q \in c$ the outcome variable of $q$ in $c$ is $v_q$ itself, so that an attribute shared by two contexts is represented by one variable and $v_q$ does not depend on which compatible attributes are measured alongside it;
\item[\textbf{D3}] (\emph{Uniformity}) one translation rule applies to every constructive blank, with no blank selectively left open;
\item[\textbf{D4}] (\emph{Marginal and verdict preservation}) for every context $c$ the joint law of $(v_q)_{q \in c}$ under $\mu$ is $P(\,\cdot \mid c)$, and the exclusive-fragment verdicts are reproduced;
\item[\textbf{D5}] (\emph{Measurement independence}) the distribution of the hidden point is independent of the settings, $P(\lambda \mid x, y) = P(\lambda)$ in the two-party scenario; this is a condition on the measure, separate from D2's context independence of the response variables.
\end{itemize}
Such a deflation is \emph{admissible}, and $\Delta$ is the class of admissible deflations; a deflation failing D2 is contextual and lies outside it. A dependence of the \emph{distribution} on the state, $\mu_s$ varying with what $s$ has learned, is the classical ignorance picture D1 provides; a dependence of the \emph{value} on the state is what the deflation is to remove, and no occurrence of an information-state parameter in a rival account is by itself a concession of the present ontology.
\end{definition}

\begin{corollary}[Common-space completion]\label{cor:bell}
In a $2 \times 2 \times 2$ Bell scenario---two parties, settings $x, y \in \{0,1\}$, outcomes $A_0, A_1, B_0, B_1 \in \{-1, +1\}$---any $\delta \in \Delta$ induces, by D1, D2 and D5, a joint law $P(A_0, A_1, B_0, B_1)$ of which the observed pairwise correlations are the marginals, by D4. By Fine's theorem such a joint law exists if and only if the complete CHSH family holds with compatible marginals \cite{fine1982hidden}; so wherever the observed correlations violate CHSH, $\Delta$ is empty. In scenarios beyond $2 \times 2 \times 2$ the corresponding marginal-polytope inequalities replace CHSH \cite{pitowsky1994george}, and the corollary is stated only where they are known.
\end{corollary}
\begin{proof}
The uniform $\delta$ acts entry-wise, and the entangled arrangement's states carry the outcome attributes of both wings under both settings as open entries, so by D1 and D2 the deflated image values all four jointly on $\Lambda$; D5 lets one measure serve all four setting pairs; D4 makes the observed correlations its marginals; and the equivalence between the existence of such a joint distribution and the Bell-type inequalities is Fine's \cite{fine1982hidden, pitowsky1994george}.
\end{proof}

\noindent The obstruction is Fine's, and the corollary applies it to the class $\Delta$: it excludes that class, not determinacy's one-place form as such, and it does not select information-state indexing from the alternatives that also reject the class. (The companion states the same exclusion as a reductio on a \emph{lossless deflation} defined by the corresponding conditions \cite{Aethus}.)

\subsection{How the two results relate}\label{sec:relation}

A deflation does one of two things with a constructive entry: it retains the entry's openness and adds a value at the same index, or it translates the openness into ignorance over a state-independent value. The dichotomy is exhaustive, since a value is either content indexed to a state or background predicated of none; the lemma meets the first and the corollary the second. The two have different force. The lemma is framework-internal, a consequence of L3 that a critic who has not granted L3 may decline; the corollary is an external, empirical obstruction, needing no ontology beyond D1--D5, and it is where the framework meets a no-go it did not write. Together they say that, on the ledger and within $\Delta$, the relabelling $\delta$ has no consistent, empirically adequate action on the framework's open entries; they do not say that the two-place reading is the only one left, which is claim (iii), argued comparatively in \S\ref{sec:placement}.

\subsection{Scope}\label{sec:scope}

Four boundaries. First, the target is the \emph{non-contextual} reparse: contextual alternatives lie outside $\Delta$, the argument does not address them, and a value $v(A, C, \lambda)$ fixed once the attribute, the context and the point are fixed may be an objective determinate fact, so contextual theories are not two-place in this paper's sense; the Kochen--Specker theorem closes richer noncontextual value assignments on independent grounds and is not leaned on \cite{kochen1967problem}. Second, Bohmian mechanics is not a counterexample but an alternative with a stated cost, compared in \S\ref{sec:hatch}. Third, the results are conditional as the introduction's table says: deny L3 and the lemma does not reach the denier, at the cost of the constructive blank the companion physics puts to work; deny L4 and both results stand, the lemma being stated at a single index and the corollary owing nothing to permanence, at the cost of historical evaluation. Fourth, D5 is a premise, and its denial is a genuine option: the superdeterministic programs keep one-place values by letting the distribution over $\Lambda$ vary with the setting choices \cite{thooft2016cellular, hossenfelder2020rethinking}. Hossenfelder and Palmer argue that violating statistical independence does not by itself cost explanatory power and that the cost must be assessed on a concrete model; the criticism entered here is the paper's own, that the correlations between hidden values and free-variable settings must be written into the initial conditions or the dynamics, and that the explanatory autonomy of experimental science is what a model which writes them there gives up. The paper declines the option on that ground; nothing in the corollary closes it. The exclusive fragment, finally, translates without incident: entries with an exclusivity clause admit the classical representation, and a reparse that stops at them reads the framework's classical fragment rather than deflating the framework, which is \S\ref{sec:hatch}'s business.

\section{Selective reparses, and what each costs}\label{sec:hatch}

A reparse that accepts it cannot deflate \emph{everything} may still hope to deflate \emph{something}: keep two-place openness where the theorems bite, restore one-place determinacy where they do not, and call the framework a local courtesy rather than a category. Every such selective reparse must say where the determinate line falls, which attributes possess one-place values and in virtue of what. The strategies actually on offer are compared here, each with the cost it states for itself.

\begin{enumerate}
\item \emph{No line: reparse everywhere.} The universal deflation of \S\ref{sec:samplespace}, closed by Corollary~\ref{cor:bell} wherever the statistics violate the Bell bounds \cite{fine1982hidden, pitowsky1994george}.
\item \emph{A preferred beable: positions, always.} The line drawn by election of an attribute family, constant as in Bohm's positions or rule-given as in the modal interpretations' spectral prescriptions; not closed, its costs a nonlocal guidance dynamics, a preferred structure that special relativity must be made to tolerate, and an equilibrium hypothesis carrying the statistics \cite{bohm1952suggested}. It is a rival ontology and not a deflation: the Bohmian does not claim the framework's blanks were her particle positions all along.
\item \emph{A line drawn by record status.} Determinacy conferred when records of an attribute proliferate irreversibly into the environment. The line-drawer concedes that determinacy is conferred by an informational circumstance, which is the framework's boundary thesis, and does not concede the second relatum's identity, her relatum being the world's record structure and not any information state, so that a recorded-but-unlearned attribute is for her one-place determinate. In the record-rich regime the two readings are empirically identical; they differ at one place only, where record supply is contested in principle, the nested-observer case of \S\ref{sec:grid}, where a record-status line without in-principle irretrievability is the for-all-practical-purposes reading and with it is the framework's own criterion. The parsimony comparison between the two bookkeepings is contested and is entered as the framework's side of a dispute, not as a settled count.
\item \emph{A line drawn by physical scale.} Mass, particle number or another dynamical threshold at which superposition gives way: a rival physics with new parameters and experimental tails, the collapse-theory row \cite{ghirardi1986unified}, testable and not a reading of anything.
\end{enumerate}

The list is a map of the strategies on offer and not a logical partition of every conceivable one. What is claimed is that each strategy actually on offer falls under one heading and carries that heading's stated cost, and that none of the costs is ``the framework, relabelled''; a selective reparse with a criterion not on the list would need its own comparison.

\section{Applications and their additional assumptions}\label{sec:grid}

This section is written to be read alone, and its claim should be stated before a single case is run, because the claim is about the \emph{list}, not about any member of it. The mechanism is one move, fixed here and never adjusted below: \emph{determinacy is a two-place relation; an open entry is a constructive, tense-free description relative to its state; resolution is transit to a successor; and apparent paradox is diagnosed by asking, of each paradoxical sentence, which index it consumed without declaring.} The applications also consume the further premises listed in the table's fourth column, each named where it enters. And one dictum organizes the counterfactual half of the yield: Peres was right that unperformed experiments have no results \cite{peres1978unperformed}, and the framework keeps his insight while distributing the modality it compressed --- in the scenario fragment, for an unperformed measurement whose alternatives are pockets of an exclusive residue, \emph{the would-counterfactual is false} whenever the residue contains a visible pocket in which the stated result fails, which it does whenever more than one outcome pocket is visible, and true only in the degenerate case in which every visible pocket agrees, an eigenstate measured in its own eigenbasis; \emph{the might-counterfactuals are true}; and \emph{the graded verdict is defined} wherever the weighting and its normalization are; agreeing superpositions are evaluated by the full framework's separate clause. No results, with truth conditions: that is the shape of every application below.

Before the cases, the schema they all instantiate, so that what follows reads as one correction applied under eight different sets of additional premises, the grade of each result read off its own row.

\begin{table}[htbp]
\centering
\caption{One mechanism, eight applications, eight sets of additional premises. Each row asks the same thing: which index did the paradoxical sentence consume without declaring? The fourth column names what each application uses beyond the correction itself.}
\label{tab:schema}
\small
\begin{tabular}{p{0.12\textwidth} p{0.2\textwidth} p{0.23\textwidth} p{0.2\textwidth} p{0.13\textwidth}}
\toprule
\textbf{Case} & \textbf{The illicit monadic question} & \textbf{The indexed question that replaces it} & \textbf{Further premises} & \textbf{Grade} \\
\midrule
Delayed choice & What did the particle really do? & What is resolved relative to which state? & Which records exist, and when & Full \\
Bomb tester & Did the photon interact with what did not happen? & Which joint descriptions survive the arrangement's record constraints, relative to which index? & The response law; persistence of the explosion and intactness records; the transparent dud & Interpretive placement; exclusion conditional \\
Cat & Is the cat really alive or dead? & What does the cat attribute resolve to relative to this state? & Permanence of the box's records (the third postulate) & Full \\
Wigner's friend & Who has the fact? & Relative to whose state is the fact determinate? & Record--coherence incompatibility; the coherent state assignment excluded only on the no-safe-refinement premise & Reconciliation full; state exclusion conditional \\
Frauchiger--Renner & Can all agents jointly apply their reasoning? & Which laboratory-state assignment is admissible relative to which index? & $P_3$, the record--coherence incompatibility, and the no-safe-refinement premise & Reconciliation full; state exclusion conditional \\
EPR & Which foliation carries the collapse? & Which state changes its indexing? & Quantitative realization imported from the physics companion & Full semantic; realization imported \\
PBR & Which underlying ontic state does $\psi$ describe? & Does the unindexed ontic-state category exist at all? & None; the row is interpretive & Interpretive; no representational claim \\
Mixtures & Is this mixture proper or improper? & Does the state assert a fact about which branch? & The exclusive/agreeing distinction & Full \\
\bottomrule
\end{tabular}
\end{table}

\paragraph{The central application: one criterion, a detector and a friend.} Two of the cases share one formal structure, and it is the proposal's clearest application, so it is set out before the cases are run. Let $A_{\mathrm{coh}}$ be a candidate state asserting a maintained coherence at a stated system boundary, not itself a base contradiction. Lemma~\ref{lem:record-exclusion}, with its converse, excludes it exactly when
\[
\forall B \le A_{\mathrm{coh}}\ \exists C \le B:\ C \in \mathrm{Base},
\]
in words, when every admissible refinement of the coherent state still has a continuation that reaches a contradiction. The quantifiers range over the full declared refinement structure, every proper descendant at every depth, and not over the present sealing of a laboratory or whether anyone has yet read a detector. Each application must therefore specify three things: the candidate state with its boundary; the stated content incompatible with the coherence, which supplies the base contradiction; and why every refinement retains a continuation reaching it.

\emph{The detector, in the specified model.} The candidate must be named exactly, because a which-path record bears on two different coherences. The detector application excludes maintained particle-path interference: an agreeing superposition of the which-path attribute in the particle's reduced description. The incompatible content is the detector's record of the path, stated jointly with that claim; here it happens to be a which-path record, but any degree of freedom carrying an outcome-distinguishing correlation would serve, and a detector's presence without a distinguishing record supplies no contradiction. The continuation is supplied by the model: the record is part of the arrangement's specified content, so every declared refinement of the coherent state has a descendant stating it, and maintained particle-path interference is excluded whether or not the record is ever read. Global particle--detector coherence is a different candidate. The entangled particle--marker state of the mixtures entry (Appendix~\ref{app:applications}) carries an orthogonal marker record and global coherence at once, its reduced particle state being mixed, and excluding that composite coherence would need its own no-safe-refinement premise for the composite, which the detector model does not supply.

\emph{The friend's laboratory, as a whole.} The candidate state maintains an agreeing superposition of the friend's outcomes, the boundary now the laboratory, and the application must supply three things that the detector model supplied for free.

\emph{The coherence claim.} The composite is the laboratory $L$, the friend together with the apparatus and whatever memory the friend keeps. Two times are kept apart. The \emph{referent} $\sigma$ is the stage the attribute describes: the outcome branches of $L$ at a stage after the friend's measurement, written $Q_\sigma$. The \emph{record time} $\tau$ is the stage at which information about $Q_\sigma$ is carried outside the boundary or acquired by an index, and nothing requires $\tau = \sigma$. The candidate asserts, relative to an outside index, the agreeing regime for $Q_\sigma$.

\emph{The base incompatibility.} An internal record does not supply it: under the ordinary quantum description a composite can carry correlated record branches and remain globally coherent, and the framework agrees, since a record inside the boundary leaves the alternatives jointly open at the boundary. The joint-invalidity relation responsible is the one that assigns openness its regime \cite{AethusUnion, Aethus}: alternatives whose distinguishing information is accessible to an index stand in the exclusive regime relative to that index. The jointly invalid pair is therefore \emph{the agreeing entry for $Q_\sigma$} together with \emph{a stated separating record of $Q_\sigma$ accessible to the same index}, and the relation is between entries with the same referent $\sigma$, whatever their record times. A later record that concerns a later attribute---the configuration at $\tau$ after an interaction that ends the coherence---has a different referent and contradicts nothing about $\sigma$. A later record that concerns $Q_\sigma$ itself does: a nondemolition copy at $\tau > \sigma$ of a memory that has persisted since $\sigma$ is an entry about $Q_\sigma$, and relative to an index carrying it the alternatives at $\sigma$ are exclusive. That is the framework's commitment, and it is where it departs from the tensed reading on which the composite ``was coherent at $\sigma$'' whatever is later learned: tense-free blankness (L2) reads the regime of $Q_\sigma$ relative to an index off what that index can access, with no date attached. Two cases thus fall under one relation: a separating record that already existed outside the boundary at $\sigma$ and is later acquired, and separating content that becomes accessible later about the attribute referring to $\sigma$.

\emph{The descent premise.} This is where the physics enters: that every admissible refinement of the coherent laboratory state retains a descendant stating a separating record of $Q_\sigma$, at $\sigma$ or at any later record time, because an observer-grade laboratory exports outcome-correlated degrees of freedom across any boundary the control can enclose, or has them copied out before the control closes. The premise does not say that the stage $\sigma$ contains an interaction across the boundary; it says that the ideal in which the alternatives at $\sigma$ are agreeing and no refinement ever carries a record of $Q_\sigma$ is unavailable for such laboratories, the record being obtained, if not already exported, later and about that stage. The paper states this as a premise about laboratories, and a physical argument would have to establish it.

\emph{Two limits.} The criterion is one, while establishing its premise is application-specific. And the conclusion concerns $A_{\mathrm{coh}}$ at the stated boundary: excluding the maintained coherence of $L$ does not by itself exclude the coherence of a larger composite, which would need its own candidate state, boundary and record claim.

\emph{A worked model of the two records.} Let the measured system $S$, an internal memory $M$ and an external degree of freedom $E$ be two-level, with the controlled boundary enclosing $S$ and $M$ and not $E$. With only the internal record, $|\Psi\rangle_{SM} = (|0\rangle_S|0\rangle_M + |1\rangle_S|1\rangle_M)/\sqrt2$ with $E$ uncorrelated, the outcome alternatives are jointly open at the boundary, no record outside it distinguishes them, the agreeing description is admissible there, and an outside measurement in the basis $|\pm\rangle = (|00\rangle \pm |11\rangle)/\sqrt2$ gives $|+\rangle$ with probability one. With a separating record, $|\Psi\rangle_{SME} = (|0\rangle_S|0\rangle_M|e_0\rangle + |1\rangle_S|1\rangle_M|e_1\rangle)/\sqrt2$ with $\langle e_0 | e_1\rangle = 0$, the distinguishing information is carried outside the boundary: relative to an index carrying the $E$-record the alternatives are exclusive, the agreeing description at the boundary of $SM$ is a base contradiction jointly with that record, the reduced state of $SM$ is the mixture $(|00\rangle\langle 00| + |11\rangle\langle 11|)/2$, and the same measurement gives $|+\rangle$ with probability one half. The record time is immaterial to the relation: if the $E$-record is made at $\tau > \sigma$ by a nondemolition copy of the persisting memory, $\mathrm{CNOT}_{M \to E}$ taking the first state to the second, it is a record of $Q_\sigma$, and the framework's description of $Q_\sigma$ relative to an index carrying it is the exclusive one. The unitary model itself says only that $SM$ was coherent before $\tau$ and mixed after; that the later copy bears on the description of $\sigma$ is the indexed reading, the premise's content and not a consequence of unitary quantum mechanics. The agreeing description is admissible again at the boundary of $SME$, which the control would have to enclose. The laboratory premise is the claim that for an observer-grade $L$ the second situation is the only one every admissible refinement reaches; the model shows what the premise asserts and does not establish it, its reduced-state calculations being standard.


\noindent Each case separates three things: the semantic clarification the correction supplies; any physical commitment, one being staked, the laboratory premise for observer-containing laboratories, under which the coherent laboratory-state assignment the Frauchiger--Renner protocol requires is inadmissible, a commitment on which the framework differs from universal unitarity and not a consequence of the indexing; and any premise about realizable records, named in the table's fourth column. The cat and the friend fall under the rule the worked model displays: a record inside the boundary leaves joint openness at the boundary and makes the alternatives exclusive relative to any index that has acquired the record. The permanence of L4 is semantic and does not imply that a physical memory cannot be coherently reversed. Throughout, what the criterion excludes is a coherent state assignment; a claim that the corresponding apparatus operation cannot be performed, or that the experimental statistics are modified, would need a further construction, left open. Each case closes with a grade on a scale declared once: \emph{full} (the paradoxical sentence consumed an undeclared index, and re-indexing leaves no residue); \emph{conditional} (the diagnosis stands modulo a named premise or open construction); \emph{interpretive} (the framework places the case without a representational or physical claim).

\paragraph{Delayed choice, and the quantum eraser.} The paradox: the experimenter's late decision --- to read or to erase the which-path marker --- appears to reach backward and settle what the particle \emph{was doing} at the slits, long after the doing. The diagnosis is one clause of the ledger: by L2, an open entry is tense-free --- it was never dated by the system's proper time, so there was no moment at the slits by which the path had to have become a fact awaiting later reversal. The path attribute stands open relative to every state lacking a record of it, however late the clock runs; the ``choice'' does not reach backward because there is nothing back there with a date on it to reach. What the delayed-choice arrangement actually exhibits is a boundary-value fact misread as an influence: the experiment's full record-structure, taken as a whole, fixes what stands resolved relative to what --- and a constraint on the whole arrangement is not a signal within it. The eraser adds bookkeeping that must be kept consistent with L4, since ``erasure'' can be heard as overwrite from the opposite direction. Three situations are distinct. A record was never acquired by a given state; a physical record existed but access to it was subsequently destroyed; a previously stated entry in a state becomes open again. The framework holds that the third never happens to the \emph{same} state: a state that carried the marker's record carries it permanently at its own index, true there forever, by L4. What the eraser accomplishes must be stated with its physical operation. Losing access to the marker---discarding it, or any operation on it that is not conditioned on an outcome---leaves the path's reduced state an equal mixture, so a later state downstream of it has the path attribute open in the exclusive, disagreeing sense, and no interference appears. Erasure proper is conditional: measuring the marker in a basis that does not distinguish the paths, and sorting the runs by that outcome, yields subensembles in which the path attribute stands open in the agreeing sense, each showing a fringe, and the fringes are complementary, so their sum shows none---provided no remaining unconditioned separating record, a copy of the marker's information carried elsewhere, suppresses the subensembles' coherence; where such a copy exists the sorted subensembles show no fringes, the path attribute standing open in the exclusive sense in each. Relative to a later state indexed to one erasure outcome, the path is agreeingly open; relative to a state indexed to none, it is exclusively open. The predecessor that resolved it stands unchanged. So the eraser changes which states resolve the attribute, never what any state resolved: re-indexing, not retro-causation, and not un-recording either. The case is the cleanest in the set because its air of paradox comes from a single undeclared tense. \emph{Grade: full.}

\paragraph{Wigner's friend.} The paradox: the friend has measured and holds an outcome; Wigner assigns the laboratory a superposition; the two descriptions seem to contradict. The diagnosis: they are two indexings of one arrangement. Relative to the friend's state the outcome attribute is resolved, a record existing there; relative to Wigner's it stands open, no record having crossed to him, and his description is correct at that address, where by L3 it is a description and not an agnosticism. The measured system's alternatives admit the exclusive, record-relative description, the friend's record carrying an outcome-distinguishing correlation, which the third postulate reads as it reads a which-path detector's; whether the whole laboratory's outcome branches can retain coherence is the separate question governed by the laboratory premise below. Indexing reconciles the two descriptions, and when Wigner opens the door, transit, not retroactive correction, brings his successor state to state what the friend's long had, both predecessors standing true of their indices. Whether the coherent state assignment for the whole laboratory across the friend's record basis is admissible is a separate matter, decided by the central application's laboratory premise and not by the permanence of a description: a permanently available description of an earlier measurement and a later reversal of its physical implementation are different objects, and excluding the assignment is not excluding the operation. \emph{Grade: reconciliation of the two descriptions, full; exclusion of the coherent state assignment for the whole laboratory, conditional on the laboratory premise; nothing claimed about the availability of the corresponding operation.}

\paragraph{Frauchiger--Renner.} The extended Wigner's-friend protocol needs a different treatment from the two-indexing case \cite{frauchiger2018quantum}. The third postulate's test has the quantifier shape that decides it: a candidate agreeing superposition is valid when some declared refinement has no contradiction-reaching descendant, and is excluded only when every permitted refinement still has one, so that neither a possible later record nor an observer's capacity to communicate excludes anything by itself. Three statements differ, and the exclusion needs the third: that the friend can communicate the outcome; that the specified experiment retains an avenue for later communication; and that every permitted refinement retains some avenue to a separating record at the relevant stage, so that no safe coherent reversal exists. Being human makes the second extremely plausible in ordinary experiments, and the passage from human record-bearing to the third is the laboratory premise of the central application, a physical claim stated as a premise and distinct from the permanence of an indexed description, neither entailing the other unaided. Wigner's friend needs only two indexings of a completed measurement; the Frauchiger--Renner argument needs an \emph{observer}, reasoning and record-keeping, to be held in agreeing superposition relative to another observer, so that one agent can treat a colleague's completed measurement as a manipulable superposition while importing that colleague's inferences.

\textit{The step the account rejects.} The original argument runs on three assumptions: (Q), that an agent who assigns a state and computes a Born probability of one may be certain of the outcome; (C), that certainty transfers between agents, an agent certain that another agent is certain of an outcome being certain of it herself; and (S), single outcomes. The laboratory premise bears on (Q), at one kind of instance: the application of (Q) by an agent outside a laboratory to a coherent state of that laboratory across its friend's recorded outcome alternatives. In the protocol that instance is used twice, by the first superobserver's friend inferring, from the coin outcome, the certain result of the second superobserver's measurement, and by the first superobserver inferring, from her own result, the spin result inside the other laboratory; both inferences are drawn from the coherent joint assignment, and under the premise that assignment is inadmissible relative to the inferring agent's index, so the certainties do not arise. The transfer rule (C) is not restricted by this paper; what the premise removes is a certainty that (C) would otherwise have transferred. Decoherence enters only as the physical process whose record consequences fix which regime the alternatives take, agreeing and exclusive superposition being defined before any weight or dynamics \cite{Aethus}.

\textit{What is excluded, and on what premise.} The candidate is the coherent description of the whole laboratory across the friend's outcome alternatives, the composite scope at which the protocol measures entire laboratories. Its exclusion requires the laboratory premise, that every declared refinement retains a descendant stating a separating record at the relevant stage; given it, $P_3$ excludes the description for the same formal reason as in the detector case, and whether actual observer-containing laboratories satisfy the premise remains open. The result is a conditional exclusion of the required coherent state assignment. It does not show that a superposition-basis measurement is unavailable, since a mixture can be measured in such a basis, with different statistics. For the ideal protocol state $|\Psi_{\mathrm{FR}}\rangle = (|h,\downarrow\rangle + |t,\downarrow\rangle + |t,\uparrow\rangle)/\sqrt3$, with $h, t$ the coin record in the first laboratory and $\downarrow, \uparrow$ the spin record in the second, the $(\mathrm{ok}, \mathrm{ok})$ cell has probability $\tfrac{1}{12}$ under the coherent assignment, and $\tfrac14$ under its replacement by the state fully dephased in the product record basis, $\rho_{\mathrm{mix}} = \tfrac13\big(|h,\downarrow\rangle\langle h,\downarrow| + |t,\downarrow\rangle\langle t,\downarrow| + |t,\uparrow\rangle\langle t,\uparrow|\big)$; dephasing the record basis of either laboratory alone leaves the cell at $\tfrac{1}{12}$, so the figure belongs to the named replacement and not to ``the mixture.'' These numbers are consequences of a particular replacement state, which the physics companion would have to supply; $P_3$ itself establishes only the conditional exclusion. \emph{Grade: reconciliation of the indexed descriptions, full; state exclusion, conditional on the laboratory premise; operation exclusion and modified statistics not claimed.}

\paragraph{The remaining applications.} The other rows of Table~\ref{tab:schema} apply the same correction under their own further premises; Appendix~\ref{app:applications} works each in full, and the bomb tester is placed in Section~\ref{sec:placement}.

\section{Objections}\label{sec:objections}

\noindent Five objections deserve answers, and each answer states what the paper establishes and what it leaves open.

\textit{Indexed determinacy is ordinary epistemic uncertainty.} This is the deflationist's objection, given its strongest form in \S\ref{sec:deflation}. The answer is the two results of \S\ref{sec:nogo}: a translation that retains the constructive openness of agreeing superpositions and adds a value at the same index contradicts the entry (Lemma~\ref{lem:preservation}), and a translation that reads the openness as ignorance over state-independent values is a single-sample-space model, which the quantum statistics exclude (Corollary~\ref{cor:bell}). The answer has the limits the paper states: where openness is exclusive, as in the sealed coin, nothing observable distinguishes the readings, and a contextual reparse is not addressed by either result.

\textit{Cross-perspective consistency.} If Wigner's state leaves open what the friend's states, do the two contradict? They do not, because they are claims with different relata: statedness relative to one state and openness relative to another are compatible, as simultaneity relative to two frames is. Consistency is required elsewhere, and what the framework guarantees is stated as a proposition, with what it does not guarantee beside it. Two states \emph{communicate} when each transits to a common refinement, a state that is a declared descendant of both; their \emph{joint assertion} $A \cdot B$ is the state whose stated content is the deductive closure of the union of their contents, a child of each.

\begin{proposition}[Agreement on shared records]\label{prop:agreement}
Let $A$ and $B$ be valid states. Assume that stated content persists under refinement, every descendant of a state stating what that state states, and that a state stating two incompatible values of one attribute is a base contradiction. Then (a) if $A$ and $B$ both state an attribute $q$, the joint assertion $A \cdot B$ lies outside the base only if they state the same value of $q$; (b) every valid common refinement of $A$ and $B$ states every attribute either of them states, with that state's value, so that a valid common refinement exists only if $A$ and $B$ agree on every attribute both state; and (c) where $A$ states $q$ and $B$ leaves it open, every valid common refinement states $q$ with $A$'s value.
\end{proposition}
\begin{proof}
(a) If the values differ, the union of the contents states both, so $A \cdot B \in \mathrm{Base}$. (b) A common refinement is a descendant of each, so by persistence it states both contents; if the values of an attribute both state differed, it would be a base contradiction, which validity excludes. (c) As in (b), $A$'s entry persisting and $B$ contributing nothing that contradicts it.
\end{proof}

\noindent A disagreement between perspectives therefore shows up as an invalid joint assertion, a definite fact about the two states together, and never as one attribute both stated and open at a single index. What the proposition does not say is that every physically admitted communication process between two observers realizes a valid common refinement of their states, or that it leaves each record intact: that is a claim about the physical realization of transit, it is not proved here or in the companion, and where an application needs it the claim is a premise of the same kind as the laboratory premise.

\textit{Retention of earlier states.} The proposal keeps earlier information states as objects, and the objection is that this is either psychology or an extravagant ontology. It is not psychology: historical evaluation, what was settled then, needs a state to be evaluated at, and the commitment is to retaining such states as semantic objects. Two things are not claimed: that a retained description entails the physical irreversibility of the records it describes, and that the permanence of a description excludes a later physical reversal. The friend case keeps the two apart.

\textit{The status of global coherence.} Does the proposal deny that a laboratory, or the universe, can stand in coherent superposition? It does not. The record-exclusion criterion excludes a specified candidate state at a specified boundary, given the incompatible record claim and the continuations the application supplies. Whether a larger system's composite coherence is excluded is a separate application, with its own candidate state and physical premise, and the paper does not run it.

\textit{Why prefer this proposal to its neighbors?} Relational quantum mechanics, QBism and Spekkens's epistemic program also make descriptions observer-relative; the proposal differs from each at identifiable points, set out in \S\ref{sec:placement}, and whether those differences are improvements is the comparative claim (iii), for the field to judge.

\section{Placement}\label{sec:placement}

Relativization of facts is occupied ground, and the map of the neighborhood is part of the position. The differentia this paper claims is narrow. The relativization is elected here, at L1, exactly as the neighbors elect theirs, and no derivational daylight is claimed on the statement itself; what \S\ref{sec:nogo} derives is an obstruction to the noncontextual, measurement-independent completion of the ledger, Fine's obstruction in this vocabulary, which the neighbors' relativizations pass as well, so it is no differentia against them. The differentia is the indexed semantics itself, proposed, together with the identification of which translations of that semantics entail the prohibited completion. That difference only functions as a difference if drawn explicitly, neighbor by neighbor.

\emph{Relational quantum mechanics} is the nearest physics neighbor: facts relative to systems, no absolute state of affairs, the relativization total \cite{rovelli1996relational}. Both programs postulate the relativization, and both pass the same obstruction to the noncontextual, measurement-independent completion, Fine's, so the positions differ not in a no-go one possesses and the other lacks but in the indexed semantics each proposes and in its consequences; Rovelli's original proposal likewise diagnoses observer-independent state descriptions as the problematic starting point. The point of contact is cross-perspective agreement. In the amendment of Adlam and Rovelli, a postulate of \emph{cross-perspective links} replaces the relativity of comparisons: when an observer has measured a variable of a system and has not since undergone an interaction destroying the information about it stored in her physical variables, a second observer who measures the physical variable carrying that information obtains a result matching hers \cite{adlamrovelli2022}. The corresponding statement here is Proposition~\ref{prop:agreement}(b)--(c): for whatever valid common refinement the two states reach, agreement on the shared attribute is a consequence of persistence and of the base, under the hypotheses stated there, and the condition that the first observer's information not be destroyed corresponds to the persistence of stated content under refinement, which the semantics has by L4 and which a physical realization of transit must respect. The difference is the location of the postulate. Relational quantum mechanics postulates the matching of results at the level of measurement outcomes; here the agreement is derived for valid joint extensions, and the physical claim, that communication realizes such an extension, is left open and marked as such in \S\ref{sec:objections}. Neither account derives the other's postulate, and no further advantage is claimed. Owed to RQM, so that the comparison does not flatter by omission: Rovelli's program never claimed a forcing result, and its not having one is a difference in ambition, not a failure of execution.

\emph{Brukner's no-go} and the Local Friendliness theorem are neighbors the contrast must keep distinct \cite{brukner2018nogo, bong2020strong}. Brukner shows that the conjunction of the universal validity of quantum theory, locality, freedom of choice and observer-independent facts yields a Bell-type inequality that quantum predictions violate; the result obstructs the conjunction, and it does not by itself show that observer-independent facts must be abandoned. Bong and colleagues derive Local Friendliness inequalities in a theory-independent way from No-Superdeterminism, Locality and Absoluteness of Observed Events; quantum theory supplies violating correlations when the superobserver has coherent control of the friend's laboratory, and their demonstration uses a photonic path as the friend's observation. Again a conjunction is obstructed. The present paper's relation to both is this. Its internal argument is a preservation result stated from the semantic ledger, while its external argument consumes the Bell-violating statistics as theirs do, so the distinction is not semantic premises against quantum premises but a preservation constraint beside a shared empirical obstruction. And where the quantum derivation of a Local Friendliness violation uses a coherent state assignment for the friend's laboratory, the laboratory premise of \S\ref{sec:grid}, if true, excludes that assignment for observer-grade laboratories while leaving the photonic demonstration untouched; whether the corresponding operations are unavailable or the observed statistics modified would need the further construction the paper leaves open. That is a point of contact between the two, stated conditionally, not a derivation of either result from the other.

\emph{QBism} shares the informational center and differs at the ontological difference \cite{fuchs2014introduction}. For the QBist, the quantum state is an agent's betting book, and measurement outcomes are experiences of the agent; the view is proudly first-personal. The present account's information states are logical objects --- individuated by stated content, truth-apt, indexable by anyone --- and what stands determinate relative to them is a structural fact about the world, not a norm on anyone's gambles. The two views will often issue the same sentence about a laboratory; they mean different things by the copula, since what stands determinate relative to an information state here is not an agent's expectation but a fact about the world's structure.

Spekkens's own program is placed here as the other candidate answer to the shared diagnosis. The toy theory reproduces a large tranche of quantum phenomenology from an epistemic restriction over underlying ontic states, an elementary system having four such states and an epistemic state representing restricted knowledge of them \cite{spekkens2007evidence}. The contrast with the present account is therefore not that one has an ontology and the other lacks one. It is between knowledge restricted over underlying ontic alternatives, one of which obtains, and constructive openness, which posits no unindexed fact beneath the blank and carries its two regimes as content. Where the toy theory reproduces the phenomena, the present account expects it to, that being the arena in which the structure of the second relatum, which states record what, dominates; where it does not, the present account offers no diagnosis of its own, and the toy theory's boundary is common data for the shared diagnosis, which answer fits it being for the comparison to decide.

\textit{A note on Einstein.} Spekkens fills the mistaken category with the quantum state taken as ontic rather than epistemic, and this paper with determinacy taken as monadic; the position here is accordingly not $\psi$-epistemic, the openness a state carries being content about the world rather than restricted knowledge of an underlying ontic configuration. The relation is the one a certain kind of realist has held toward Einstein: the diagnosis of incompleteness taken as load-bearing, the particular repair treated as a separable conjecture its author grades as unfinished.

\emph{Modal interpretations and consistent histories} occupy, respectively, the beable strategy of \S\ref{sec:hatch} and a sophisticated cousin of it. The modal strategy --- definite values assigned by rule to a preferred set of attributes \cite{lombardi2021modal} --- is a principled line-drawer, and \S\ref{sec:hatch} states its family's cost rather than refuting it. Consistent histories \cite{griffiths2002consistent} relativizes assertions to a framework of histories and forbids mixing frameworks --- a relativization, but to the \emph{theorist's descriptive choice}, where the present account relativizes to \emph{states in the world}: laboratory records index facts whether or not anyone selects a framework, and the single-framework rule is replaced by the indexing discipline, which no theorist administers.

\emph{Fine's fragmentalism} is the philosophical cousin, and to our knowledge the comparison has not been drawn in print. Fine's tense-realist proposal lets reality itself divide into standpoints --- fragments internally coherent, jointly irreconcilable, with no privileged global perspective \cite{fine2005tense}. The structural kinship with indexed determinacy is plain, and the present account can be put in his vocabulary as \emph{fragmentalism with an engine}: the fragments are information states rather than metaphysical postulates; the coherence conditions within and across fragments are the ledger's stated rules rather than primitive constraints; and --- the engine --- the fragmentation is not merely adopted: \S\ref{sec:nogo} shows that, once the ledger's constructive reading of blankness is granted, the fragments cannot be uniformly re-read as noncontextual ignorance over pre-existing values while the Bell statistics stand---supplying the \emph{why} that fragmentalism, as a pure metaphysics, was free to leave open, at the strength those results have. If the comparison holds up, each side gains: the framework inherits a worked philosophical genus, and fragmentalism inherits a physics in which its structure does work.

Finally, the competitors on the dissolution claim itself, since the argumentative \emph{shape} of \S\ref{sec:grid} --- one commitment, the whole paradox set --- is not unique. Zwirn's convivial solipsism dissolves the set, Wigner's friend and its sharpenings included, by denying that measurement effects any physical change of state; the view is deliberately first-personal to the edge of solipsism, and its dissolution is a reading rather than a derived result \cite{zwirn2023convivial}. A recent program of interaction-conditional semantics runs configuration-relative predication over substantially the same case list, diagnosing the paradoxes as category errors, discusses the Bell and Kochen--Specker constraints expressly, and positions itself near the contexts-systems-modalities ontology of Auff\`eves and Grangier and Healey's pragmatist realism \cite{ics2026dissolution, auffeves2016contexts, healey2012quantum}. It disclaims observer-dependence, its viewpoints being physical configurations rather than mental states, so it stands on the objectivity-preserving side of the field beside the present account, and of the dissolution rivals it is the nearest. The differences are in the positive content posited, and they are listed rather than evaluated. The present account's index is an information state, a deductively closed body of stated content, not a physical configuration. Its blank is content with two regimes, the agreeing and the exclusive, whose assignment the composition rule fixes \cite{AethusUnion}. Its counterfactuals are evaluated by an operator on that content, reopening at a declared scope \cite{AethusSemantics}. Its validity predicate is the fixed point of \S\ref{sec:ledger}, which yields the one record criterion of \S\ref{sec:grid} for a detector and a laboratory alike. And it retains predecessor states as objects, which is what claim (i) is stated over. Whether configuration-relative predication can carry these structures, retained configurations among them, is a question about that program's machinery which its text and not this paper must answer; nothing here shows that it cannot. The results of \S\ref{sec:nogo} do not show that configuration-relative predication is wrong or that monadic determinacy is impossible; they show that \emph{this} relational semantics, with blankness read constructively, cannot be uniformly re-read as noncontextual ignorance over pre-existing values while measurement independence and the Bell-violating statistics stand, and the obstruction half of that is shared by every account that rejects the same completion class. The category-mistake thesis is accordingly not what those results establish. It is the explanatory hypothesis the cross-case fit of \S\ref{sec:grid} supports, claim (iii), and the comparison with these programs is where it has to be argued.


\paragraph{Bomb tester.} The Elitzur--Vaidman arrangement \cite{elitzur1993quantum} illustrates the tense-free indexing at its sharpest. The attribute at issue, whether the bomb had the specified absorbing response at the trial event, has a fixed spacetime referent; later evidence may settle it without changing the event it describes, and an outcome's being unobserved is a fact about the index while its record structure is a fact about the arrangement. A dark click identifies a live bomb without an explosion in that run, and the framework reads it as an identifying outcome in the admissible residue of the arrangement, excluding the candidate agreeing description of two recombining amplitudes conditionally, on the persistence of the bomb's explosion and intactness records, by the same lemma that governs the detector and the friend. The dark click is not evidence that the macroscopic past must be ontologically blank for every interpretation; it is the framework's own placement, with its premises named, and the long paper carries the treatment, including the further case of a live-or-dud response held in agreeing superposition, which it decides on its definition of the would-clause and a detector premise, the engaged live response rendered permanently accessible, as a proposed divergence from universal unitarity with an explicit test; that argument is the long paper's and is only pointed to here.


\section{Limits, and what remains open}\label{sec:limits}

What follows is the list of what remains open, graded as the applications were.

\emph{The realization bridge.} This paper argued that a semantics is committed to nonclassical structure; whether that structure is physically instantiated --- whether the world's records and statistics are as the indexed picture requires --- is the physics companion's business, argued there and merely cited here \cite{AethusQuant, AethusBrief}. A reader persuaded by every section above has been given a category, a preservation lemma and an obstruction, not a completed physics. The no-go is exactly this: the specified noncontextual, measurement-independent completion cannot preserve the required statistics, and the obstruction does not uniquely select information-state indexing from all remaining interpretations.

\emph{The weights.} The graded openness that L3 grants and \S\ref{sec:grid} spends is taken as structure; its derivation --- why those weights, from what --- is an assigned successor program, open on the books, and nothing above pretends otherwise.

\emph{The irretrievability construction.} The Frauchiger--Renner verdict is conditional on it, and was graded accordingly; the permanently sealed laboratory --- unretrievability over the whole of the future --- remains the contested case, and the construction's completion or failure is the single edge along which \S\ref{sec:grid}'s hardest entry moves. One kinship should be named at that edge, plainly, and its limit with it: on this flank the position shares with the dynamical-collapse programs the \emph{position} that coherence ends somewhere for observer-grade systems, and shares nothing of their mechanism. No criterion is posited beyond the third postulate together with the records premise for the laboratory stated above, the postulate excluding a candidate when every declared refinement of it retains a contradiction-reaching descendant, Lemma~\ref{lem:record-exclusion}, the same rule that ends interference when a detector carries a distinguishing record, so the framework draws no cut of its own; what the resolution is conditional on is one physical claim, that an observer-grade record structure exports its records unblockably, an open question about the physics of records and not a parameter, which is where it differs from GRW's conditionality on a rate.
% >>> EDIT S5 (Fable 5, Aug 29 2026): the ceiling graded as a commitment to a prediction, with the two disciplines
Two constraints keep that commitment from going verbal, and both are owed here as they are owed in the overview \cite{AethusBrief}. First, ``unblockable'' must not be cashed as what unitary dynamics already calls irreversible decoherence, on which reading the two sides agree on every experiment and the claim is a gloss on a shared prediction; for it to differ in kind, the record must be one unitarity still predicts recoverable in principle --- the spin-echo and Loschmidt-echo regime, records scrambled but present --- and the construction's job is to say on which side of that regime an observer's records fall. Second, the interferometry program climbing toward ever-richer record-bearing systems held coherent can only \emph{fail to find} the point at which the escape closes: every successful recombination reads as ``records still blockable'' until the construction names the conditions. The claim therefore stands as a \emph{commitment to} a prediction rather than a prediction, and becomes the latter at exactly the point where the conditions are named; there is no rate here and no scale, only the third postulate and an open question about records.
% <<< END EDIT S5
\emph{The ledger itself.} The position adopts L1--L4; each result uses the premises identified with it in the introduction's table. Deny L3 and the constructive blank collapses to a gap, and claim (i) does not reach the denier; deny L4 and historical evaluation loses its objects, both results of \S\ref{sec:nogo} standing as stated; costing a denial is not refuting it, and an opponent who pays is a live opponent.

What would refute the paper, stated so that the criterion cuts both ways: a demonstration that the counterfactual clauses can run without persisting predecessors would remove the historical-evaluation justification for L4, leaving both results of \S\ref{sec:nogo} untouched and the ontology without that support; a demonstration that the deflated image escapes the single-sample-space characterization would void Corollary~\ref{cor:bell}; a selective reparse with a criterion not on the list of \S\ref{sec:hatch} would call for a comparison the paper has not made. These are the paper's exposed edges, and it prefers to name them than to be caught defending them by surprise.



\appendix
\section{Proper children: the coefficient conventions}\label{app:properchild}

In the framework's additive definition, $B$ is a proper child of $A$, written $B \le A$ as the order of the Aethic tree ($B \le_\beta A$ with its weight $\beta$), when some admissible declared decomposition $A \rightsquigarrow bB + cC$, not an identity of the weighted algebra, with $cC$ an admissible remainder, has $b$ nonzero and not measure-zero relative to the remainder's weight $c$; its reported weight is $\beta = b/(b + c)$ where the measure type admits division, taken at $B$'s maximum coefficient within a fixed free-sum representation, by the Relational Weight Principle (in the algebra companion's strict relative-negligibility relation, $b \ne 0$ and $b \not\ll c$). A \emph{child} is $B = A \cdot C$, multiplicatively, the joint assertion of $A$ and $C$. Both relations are transitive, childhood directly and proper childhood under the algebra companion's admissibility and coefficient-composition assumptions.

\section{The remaining applications in full}\label{app:applications}

\paragraph{Schr\"odinger's cat.} The paradox: the formalism assigns the unobserved cat a superposition of alive and dead, and common sense revolts. The diagnosis: ``the cat is in superposition'' is elliptical for a two-place claim --- the cat's vital attribute stands open \emph{relative to the outside state}, jointly, at the weights the arrangement fixes --- and so read, the sentence is true and unmysterious: it is exactly what L3 says an open entry is, predicated of a state that genuinely lacks the resolving record. The revolt of common sense is aimed at a different sentence, the one-place reading --- the cat, full stop, neither alive nor dead --- which the framework never asserts of anything. What discharges the case's macroscopic flavor is not a scale threshold, since none is consulted: it is the fact that macroscopic arrangements \emph{supply records} promiscuously --- a cat's body is a record-generating engine --- so that the states relative to which the attribute stands open are contrived and short-lived, and the ordinary situation is the resolved one. Where exactly the record supply draws the working boundary is a classification question for the physics companion \cite{AethusQuant}, deliberately not imported here; the dissolution needs only that the boundary is drawn by records, not by size, and that the paradoxical sentence was one-place all along. \emph{Grade: full.}

\paragraph{EPR, and the relativistic collapse worry.} The paradox: a measurement on one wing appears to change the distant wing's state instantaneously, and relativity asks along which foliation the change propagates. The diagnosis: the near-wing measurement re-indexes rather than acts, changing which states resolve the distant attribute, beginning with the near experimenter's successor. Three claims are kept apart: (i) no superluminal causal influence; (ii) no change in the distant system's spacetime state; (iii) no change in which facts are determinate relative to a given index, relative to Alice's post-measurement state Bob's outcome being correlated and relative to Bob's pre-measurement state open, neither state changing because the other was updated. The framework asserts all three and surrenders Bell factorizability, a fourth claim and the only one the experiments refute, already spent in \S\ref{sec:nogo}. The foliation question asks on which slice a one-place fact flipped, and there is no one-place fact to flip, so no preferred slicing is owed. What is preserved is semantic locality. The quantitative no-signalling identities, $P(a \mid x, y) = P(a \mid x)$ and $P(b \mid x, y) = P(b \mid y)$, are not derived by the semantics; they belong to the physics companion \cite{AethusQuant}, as an identity of the imported weight law. The account carries no classical causal model of the correlations, and so no tuned parameter, which is not the claim of having none to hide: an arrowless coupling model can signal, and the fine-tuning question \cite{wood_spekkens_2015} is answered, where it is, by the imported quantum restriction. \emph{Grade: full on the semantics; physical realization imported.}

\paragraph{PBR.} The theorem of Pusey, Barrett and Rudolph is standardly read as forcing the quantum state's ontic status \cite{pusey2012reality}. The framework's relation to it is non-engagement, and the claim is interpretive: the theorem quantifies over un-indexed ontic states, configurations $\lambda$ that systems have independently of any information state, the category \S\ref{sec:mistake} retires. Enlarging the variable to include the index, $\lambda' = (\lambda, \text{Aethus})$, makes it relational, a state of a system-plus-index pair, and the framework declines to read it as the system variable the ontological-models framework requires, a recognized exit \cite{leifer2014psi}. Different agents can hold different distributions over preparations, but relative to any given state the other agent is herself an exclusive open entry whose pockets hand back that state's own weighting, so two overlapping distributions never co-describe one un-indexed ontic state. \emph{What this does not show.} The refusal to read the enlarged variable as a system state is not a proof that no representation of PBR's kind exists: whether the framework's statistics admit a measurable space of ontic states, with preparations as distributions, preparation independence as a product law and the quantum statistics reproduced, is not settled here, and any representation that did would be $\psi$-ontic by the theorem. \emph{Grade: interpretive; no representational claim.}

\paragraph{Proper and improper mixtures.} The pressure point, pressed hardest by d'Espagnat: the reduced state of an entangled subsystem is mathematically a mixture, yet it cannot be read as ignorance of which pure state obtains \cite{despagnat1976conceptual}. The framework keeps each claim attached to its object. A proper mixture is an exclusive open entry: the state's content includes that exactly one alternative stands, ignorance in the ordinary sense. An improper mixture is the reduced description of a subsystem of a larger state, and the coherence it lacks is located in that larger state: the composite carries the coherence, and the subsystem has its reduced density operator. The entangled particle--marker state is the standard instance, an orthogonal marker record and global coherence at once with a mixed reduced state for the particle. Which regime the subsystem's alternatives are assigned is fixed relative to the specified record structure: the pruning law's accessibility trigger licenses the exclusive reading where a record of the distinguishing information is accessible and forbids it where that information is permanently unavailable \cite{Aethus}. ``Improper'' is therefore not a synonym for ``agreeing'': the regime belongs to the alternatives given the records, not to the reduced operator. Equal reduced density operators give identical statistics for every measurement on the subsystem alone, $\Pr(k) = \mathrm{Tr}(\rho_S E_k)$ for every POVM $\{E_k\}$, so no local measurement discriminates the mixtures; what separates them is access to the purifying or record degrees of freedom, as in the eraser configuration, whose quantitative face is the physics companion's \cite{AethusQuant}. The same record structure fixes which ensemble decomposition a proper mixture carries, a question the density operator alone cannot answer. \emph{Grade: full, relative to the specified record structure.}

\printbibliography

\end{document}